\documentclass[12pt]{article}
\usepackage{graphicx} % Required for inserting images
\usepackage[utf8]{inputenc}
\usepackage{mathtext}
\usepackage[T2A]{fontenc}
\usepackage[english, russian]{babel}
\usepackage[left=2.5cm, right=2.0cm, top=2.0cm, bottom=2.0cm]{geometry}
\usepackage{verbatim}
\usepackage{epigraph}
\usepackage{multirow}
\usepackage{amsmath, amsfonts, amssymb}


\date{}

\begin{document}
\selectlanguage{russian}
\textbf{Введение}

Любые научные исследования, касающиеся использования и обработки полученных данных в различных областях, как в социологии, экономике, физике, так и в педагогике, не могут обойтись без использования математического аппарата, поскольку только строгий математический анализ может однозначно определить эффективность проведенного эксперимента. Для этих целей применяют, как правило, статистические методы при проверке предположений.

\textbf{\textit{Методы математической статистики}} позволяют выявить различные закономерности по выборочным данным, выдвигать статистические гипотезы и делать обоснованные выводы, используя математическое обоснование. При проведении педагогического эксперимента часто возникают исследовательские задачи, связанные с применением новой методики преподавания, новой образовательной технологии. Методы проверки гипотез могут быть использованы в данном случае для подтверждения эффективности новой методики.

Сравнивая на глазок (по процентным соотношениям) результаты до и после какого-либо воздействия (эксперимента), исследователь может прийти к заключению, что если наблюдаются различия в результатах, то имеет место различие в сравниваемых выборках. Однако подобный подход категорически неприемлем, так как для процентов нельзя определить уровень достоверности в различиях. Проценты, взятые сами по себе, не дают возможности делать статистически достоверные выводы. Чтобы доказать эффективность какого-либо воздействия, необходимо выявить статистически значимую тенденцию в смещении (сдвиге) показателей. Для решения подобных задач исследователь может использовать ряд \textbf{\textit{критериев различия}}, о которых пойдет речь в данной работе.

\medskip
\textbf{Непараметрические критерии методов математической статистики}

\epigraph{Небольшие различия в начальных учловиях раждают огромный различия в конечком явлении... Предсказание становится невозможным}{\textit{Жюль Анри Пуанкаре}}

\textbf{Общие определения}

\medskip
\textbf{\textit{Статистические гипотезы}} – это предположения или допущения о неизвестных параметрах, выражаемых в терминах вероятности, которые могут быть проверены на основании выборочных показателей с помощью статистических критериев.
Статистические гипотезы различают по виду предположений, содержащихся в них, при этом рассматривают \textit{нулевую} и \textit{альтернативную} гипотезы.

\medskip
Существует два типа статистических методов:
\begin{enumerate}
    \item \textbf{Параметрические методы} – это количественные методы статистической обработки данных, применение которых требует обязательного знания закона распределения изучаемых признаков в совокупности и вычисления их основных параметров.
    \item \textbf{Непараметрические методы} – это количественные методы статистической обработки данных, применение которых не требует знания закона распределения изучаемых признаков в совокупности и вычисления их основных параметров.
\end{enumerate}
Статистический критерий (К) (критерий значимости) – это некий параметр, вычисленный по определенному алгоритму, который используется для проверки основной гипотезы. В соответствии в применяемыми методами статистические критерии также бывают:
\begin{enumerate}
    \item Параметрические критерии:
    \begin{enumerate}
        \item критерий Стьюдента;
        \item критерий Фишера,
        \item критерий Романовского.
    \end{enumerate}
    \item Непараметрические критерии:
    \begin{enumerate}
        \item критерий знаков;
        \item критерий Пирсена;
        \item критерий Манна-Уитни.
    \end{enumerate}
\end{enumerate}

\textbf{Критерий знаков (G-критерий)}

\medskip
\textbf{Описание критерия знаков}

\medskip
Данный критерий был разработан доктором \textbf{\textit{Джоном Арбетнотом}} – шотландским врачом, сатириком и эрудитом в 1710 году.

Критерий предназначен для сравнения состояния некоторого свойства у членов двух зависимых выборок на основе измерений.

Имеется две серии наблюдений над случайными переменными $X$ и $Y$, полученные при рассмотрении двух \textbf{зависимых выборок}. На их основе составлено N пар вида $(x_i, y_i)$, где $x_i, y_i$ – результаты двукратного измерения одного и того же свойства у одного и того же объекта.

Элементы каждой пары $x_i, y_i$ сравниваются между собой по величине, и паре присваивается знак «+», если $x_i < y_i$ , знак «–», если $x_i > y_i$  и «0», если $x_i = y_i$.

\textit{Нулевая гипотеза} формулируются следующим образом: в состоянии изучаемого свойства нет значимых различий при первичном и вторичном измерениях.

\textit{Альтернативная гипотеза}: законы распределения величин $X$ и $Y$ существенно различны в одной и той же совокупности при первичном и вторичном измерениях свойства.

\bigskip
\textbf{Статистика критерия} ($T_{эмп}$) определяется следующим образом:
допустим, что из $N$ пар $(x, y)$ нашлось несколько пар, в которых значения $x_i$ и $y_i$ равны. Такие пары обозначаются знаком «0» и при подсчете значения величины $T_{эмп}$ не учитываются. Предположим, что за вычетом из числа $N$ числа пар, обозначенных знаком «0», осталось всего n пар. Среди оставшихся $n$ пар подсчитаем число пар, обозначенных знаком «+». Это значение и будет определять статистику критерия $T_{эмп}$.

Нулевая гипотеза принимается на уровне значимости 0,05, если наблюдаемое значение $T_{эмп}$ < $Т_{крит}$, где значение $Т_{крит}$ определяется из статистических таблиц для критерия знаков.

\medskip
\textbf{Применение критерия знаков}

\medskip
\textbf{Пример 1.} Учащиеся выполняли контрольную работу, направленную на проверку усвоения некоторого понятия. Пятнадцати учащимся затем предложили электронное пособие, составленное с целью формирования данного понятия у учащихся с низким уровнем обучаемости. После изучения пособия учащиеся снова выполняли ту же контрольного работу, которая оценивалась по пятибалльной системе.
Результаты двукратного выполнения работы (в баллах) 15 учащимися запишем в форме таблицы (см. табл. ).

\begin{table}[h]
    \centering
    \caption{Результаты контрольных работ}
    \bigskip
    \begin{tabular}{|c|c|c|c|c|c|c|c|c|c|c|c|c|c|c|c|} \hline
         \multirow{2}{*}{}&\multicolumn{15}{c|}{\textbf{Учащиеся(№)}} \\ \cline{2-16}
        &1&2&3&4&5&6&7&8&9&10&11&12&13&14&15 \\ \hline
        Первое выполнение&2&2&2&2&2&3&3&3&4&5&2&2&3&3&3 \\ \hline
        Второе выполнение&2&3&3&4&3&2&3&4&4&3&4&3&2&4&4 \\ \hline
        Разность отметок&0&+&+&+&+&-&0&+&+&0&+&+&-&+&+ \\ \hline
    \end{tabular}
\end{table}

Более наглядно результаты эксперимента можно представить на диаграмме (см. рис. ).
 
Проверяется гипотеза $H_0$: состояние знаний учащихся не повысилось после изучения пособия. Альтернативная гипотеза $H_1$: состояние знаний учащихся повысилось после изучения пособия.

Подсчитаем значение статистики критерия $T_{эмп}$ равное числу положительных разностей отметок, полученных учащимися. Согласно данным табл. 1 $T_{эмп}=10, n=12$.

Для определения критических значений статистики критерия $Т_{крит}$ используем соответствующие таблицы. Для уровня значимости $\alpha = 0,05$ при $n = 12$ значение $Т_{крит} = 9$. Следовательно выполняется неравенство $Т_{эмп} > Т_{крит}$ (10 > 9). Поэтому в соответствии с правилом принятия решения нулевая гипотеза отклоняется на уровне значимости 0,05 и принимается альтернативная гипотеза, что позволяет сделать вывод об улучшении знаний учащихся после самостоятельного изучения пособия.

\medskip
\textbf{Критерий Пирсона (хи-квадрат)}

\medskip
\textbf{Суть критерия хи-квадрат}

\medskip
Предложенный в 1900 году английским математиком \textbf{\textit{Карлом Пирсоном}} критерий согласия $\chi^2$ (хи-квадрат) применяется для сравнения распределений объектов двух совокупностей на основе измерений по шкале наименований в двух \textbf{независимых выборках}.

\bigskip
Предположим, что состояние изучаемого свойства (например, выполнение определенного задания) измеряется у каждого объекта по шкале наименований, имеющей только две взаимоисключающие категории (например, верно – неверно). По результатам измерения состояния изучаемого свойства у объектов двух выборок составляется четырехклеточная таблица (см. табл. ).

\begin{table}[h]
    \centering
    \caption{Результаты измерения}
    \bigskip
    \begin{tabular}{c|c|c|c} 
         &Категория1&Категория2&  \\
         \cline{2-3} Выборка № 1&\textit{$O_{11}$}&\textit{$O_{12}$}&\textit{$O_{11} + O_{12} = n_1$} \\
         \cline{2-3}Выборка № 1&\textit{$O_{21}$}&\textit{$O_{22}$}&\textit{$O_{21} + O_{22} = n_2$} \\
         \cline{2-3}&\textit{$O_{11} + O_{21}$}&\textit{$O_{21} + O_{22}$}&\textit{$n_1 + n_2 = N$}  \\
    \end{tabular}
\end{table}
 
В этой таблице $О_{ij}$ – число объектов в $i-ой$ выборке, попавших в $j-ую$ категорию по состоянию изучаемого свойства; $i=1,2$ – число выборок; $j=1,2$ – число категорий; $N$ – общее число наблюдений, равное  $n_1+n_2$ или вычисляемое по формуле:

\begin{equation}
    N=\sum_{i,j=1}^{2} O_{ij}
\end{equation}

Тогда на основе данных таблицы  можно проверить нулевую гипотезу о равенстве вероятностей попадания объектов первой и второй совокупностей в первую (вторую) категорию шкалы измерения проверяемого свойства, например, гипотезу о равенстве вероятностей верного выполнения некоторого задания учащимися контрольных и экспериментальных классов.
При проверке нулевых гипотез не обязательно, чтобы значения вероятностей $p_1$ и $p_2$ были известны, так как гипотезы только устанавливают между ними некоторые соотношения (равенство, больше или меньше).

Для проверки рассмотренных выше нулевых гипотез по данным таблицы 2 подсчитывается значение статистики критерия $Т_{эмп}$ по следующей формуле:

\begin{equation}
    T=\frac{N\cdot (|O_{11} \cdot O_{22} -O_{12} \cdot O_{21}| - \frac{N}{2})^2}{n_1 \cdot n_2 \cdot (O_{11} + O_{21}) \cdot (O_{12} + O_{22})}
\end{equation}
                                            
где $n_1, n_2$ – объемы выборок, $N$ – общее число наблюдений, рассчитанное по формуле.

Проводится проверка гипотезы $H_0: p_1 \leq p_2$, при альтернативе – $Н_1: р_1 > р_2$. Пусть $\alpha$ – принятый уровень значимости. Тогда значение статистики $Т_{эмп}$, полученное на основе экспериментальных данных, сравнивается с критическим значением статистики $Т_{крит}$, которое определяется по таблице  с учетом выбранного значения $\alpha$. Если верно неравенство $T_{эпм}<Т_{крит}$, то нулевая гипотеза принимается на уровне $\alpha$. Если данное неравенство не выполняется, то у нас нет достаточных оснований для отклонения нулевой гипотезы.

\medskip
\textbf{Применение критерия хи-квадрат}

\medskip
\textbf{Пример 2.} Проводился эксперимент, направленный на выявление лучшего из учебников, написанных двумя авторскими коллективами для изучения геометрии 9 класса. Для проведения эксперимента методом случайного отбора были выбраны два района. Учащиеся первого района (20 классов) обучались по учебнику № 1, учащиеся второго района (15 классов) обучались по учебнику №2.

Рассмотрим методику сравнения ответов учителей экспериментальных школ двух районов на один из вопросов анкеты: "Доступен ли учебник в целом для самостоятельного чтения и помогает ли он усвоить материал, который учитель не объяснял в классе?" (Ответ: да – нет.)

Ответы 20 учителей первого района и 15 учителей второго района распределим на две категории и запишем в форме таблицы.

\begin{table}[h]
    \centering
    \caption{Ответы учителей}
    \bigskip
    \begin{tabular}{c|c|c|c} 
         &Да1&Нет&  \\
         \cline{2-3} Выборка № 1&\textit{$O_{11} = 15$}&\textit{$O_{12} = 5$}&\textit{$n_1 = O_{11} + O_{12} = 20$} \\
         \cline{2-3}Выборка № 1&\textit{$O_{21} = 7$}&\textit{$O_{22} = 8$}&\textit{$n_2 = O_{21} + O_{22} = 15$} \\
         \cline{2-3}&\textit{$O_{11} + O_{21} = 22$}&\textit{$O_{21} + O_{22} = 13$}&\textit{$N = n_1 + n_2 = 35$}  \\
    \end{tabular}
\end{table}

Подсчитаем статистику критерия по формуле :

$$
    T=\frac{35 \cdot (|15 \cdot 8 - 5 \cdot 7| - \frac{35}{2})^2}{20 \cdot 15 \cdot (15 + 7) \cdot (5 + 8)}
$$
По таблице для одной степени свободы ($v = l$) и уровня значимости $\alpha$ = 0,05 найдем $Т_{крит} = 3,84$. Отсюда верно неравенство $Т_{эмп} < Т_{крит}$ (1,86 < 3,84). Согласно правилу принятия решений, полученный результат не дает достаточных оснований для отклонения нулевой гипотезы, то есть результаты проведенного опроса учителей двух экспериментальных районов не дают достаточных оснований для отклонения предположения об одинаковой доступности учебников №1 и 2 для самостоятельного чтения учащимися.

\medskip
\textbf{U-критерий Манна – Уитни}

\medskip
Этот метод сравнения двух выборок признается наиболее чувствительным и мощным среди прочих непараметрических критериев. Он был предложен в 1945 году \textbf{\textit{Фрэнком Уилкоксоном}}. В 1947 году он был существенно переработан и расширен \textbf{\textit{Х.Б.Манном}} и \textbf{\textit{Д.Р.Уитни}}, по именам которых и называется.

\medskip
\textbf{Математическая основа U-критерия}

\medskip
Согласно нулевой гипотезе, сравниваемые совокупности имеют одинаковые распределения. Техника метода состоит в том, что все варианты сравниваемых совокупностей ранжируют в одном общем ряду: каждому значению присваивают ранг, порядковый номер. При этом одинаковым (повторяющимся) значениям вариант должен соответствовать один и тот же средний ранг (они как бы «делят места»). После этого ранги вариант суммируют отдельно по каждой выборке:

\begin{equation}
    R_1=\sum_{i=1}^{n_1} r_i
\end{equation}

\begin{equation}
    R_2=\sum_{j=1}^{n_2} r_j
\end{equation}

и вычисляют величину критерия:

\begin{equation}
    T =\frac{U - 0.5 \cdot n_1 \cdot n_2}{\sqrt{(n_1 \cdot n_2 \cdot (n+1)/12)}}
\end{equation}

где $U = max(U_1, U_2)$ – максимальное значение из двух величин:

\begin{equation}
    U_1 = n_1 \cdot n_2 + 0.5 \cdot n_1(n_1+1)-R_1
\end{equation}

\begin{equation}
    U_2 = n_1 \cdot n_2 + 0.5 \cdot n_2(n_2+1)-R_2
\end{equation}
                                           
Если выборка достаточно велика ($n > 20$), величина статистики $Т_{эмп}$ сравнивается с табличным значением критерия Стьюдента для $df = \infty$ и $\alpha = 0.1$ (то есть для верхней 95\verb!%! области нормального распределения).

\medskip
\textbf{Применение U-критерия}

\medskip
\textbf{Пример 3.} Сравним результаты тестирования групп $X$ и $Y$, обучаемых по различным методикам, представленные в табл. . 

\begin{table}[h]
    \centering
    \caption{Результаты тестирования}
    \bigskip
    \begin{tabular}{|c|c|c|c|c|c|c|c|c|c|} \hline
         \textit{X} &29&31&38&41&43&44&45&71&  \\ \hline
         \textit{Y} &31&50&51&52&52&52&54&62&32  \\  \hline
    \end{tabular}
\end{table}
Более наглядно результаты эксперимента можно представить на графике (см. рис. ).
 
Из таблицы  следует, что $n_1 = 8, n_2 = 9$. Высокие коэффициенты вариации (30 и 17\verb!%!) говорят о том, что распределения признаков, скорее всего, не соответствуют нормальному. Поэтому сравнивать средние следует с помощью непараметрического критерия.

Упорядочим все значения $X$ и $Y$ в новую таблицу  так, чтобы значения каждой выборки располагались в двух отдельных рядах ($X, Y$). Такое расположение упрощает назначение рангов (ряды $r_1, r_2$) и суммирование рангов ($R$), рассчитанные по формулам и .

\begin{table}[h]
    \centering
    \caption{Ранжирование результатов тестирования}
    \bigskip
    \begin{tabular}{|c|c|c|c|c|c|c|c|c|c|c|c|c|c|c|c|c|c|c|} \hline
        \textit{№} &1&2&3&4&5&6&7&8&9&10&11&12&13&14&15&16&17&R \\ \hline
        \textit{X} &29&&31&38&41&43&44&45&&&&&&&&&71&  \\ \hline
        \textit{Y} &&31&&&&&&&50&51&52&52&52&54&62&62&&  \\  \hline
        \textit{$r_1$} &1&&2.5&4&5&6&7&8&&&&&&&&&17&50.5  \\  \hline
        \textit{$r_2$} &&2.5&&&&&&&9&10&12&12&12&14&15.5&15.5&&102.5  \\  \hline
    \end{tabular}
\end{table}

Вычисляем по формулам и :
$$
    U_1 = 9 \cdot 8 + 0.5 \cdot 9 \cdot (9+1) - 50.5 = 66.5
$$
$$
    U_2 = 9 \cdot 8 + 0.5 \cdot 8 \cdot (8+1) - 102.5 = 5.5   
$$
$$
    U = max(U_1, U_2) = 66.5
$$
  
Подставляя найденные значения в формулу, находим $Т_{эмп}$:
$$
    T =\frac{66.5 - 0.5 \cdot 9 \cdot 8}{\sqrt{(9 \cdot 8 \cdot 10/12)}} = 3.81
$$
  
Полученное значение (3.81) больше табличного $Ткрит(0.1, \infty) = 1.65$, следовательно, различия между выборками достоверны.

\medskip
\textbf{Заключение}

\medskip
Для более достоверного подтверждения результатов исследований, их визуального представления используются методы математической статистики, которые могут быть как параметрическими, так и непараметрическими.
Первые обладают большей статистической мощностью в сравнении со вторыми. Это означает, что если в природе существует какая-либо закономерность, то с наибольшей вероятностью эта закономерность будет выявлена с помощью параметрических критериев. Однако такие критерии возможно применять только при одновременном соблюдении двух условий: количественном характере вариации и нормальном распределении.

Если же эти условия не выполняются (например, высокие коэффициенты вариации могут указывать на то, что распределения признаков, скорее всего, не соответствуют нормальному), то используют непараметрические критерии, которые обычно проще и не требуют сложных расчетов.

В данной работе были рассмотрены три непараметрических критерия методов математической статистики: критерий знаков, критерий Пирсона (хи-квадрат) и U-критерий Манна – Уитни. По каждому критерию приведено подробное описание алгоритма расчетов, указаны границы применимости и приведен конкретный пример применения критерия для оценки результатов эксперимента.

Следует, однако, отметить, что список непараметрических критериев не ограничивается указанными. Выделяют также следующие непараметрические критерии:

\begin{itemize}
    \item Q-критерий Розенбаума;
    \item критерий Колмогорова-Смирнова;
    \item критерий Андерсона-Дарлинга;
    \item критерий согласия Кёйпера
    \item и др.
\end{itemize}

\end{document}
